\topic{force-one-form}{Force as a one-form, and what that makes impulse}

Reading anchor: Sutton and Biblarz, ninth edition, Section 2.1, printed
p.~26 / PDF p.~50 \cite{sutton}; Frankel, Overview Section O.f, ``Line
Integrals'', printed p.~xxxix \cite{frankel}; Hill and Peterson, Sections 1.2
and 7.2, printed pp.~4 and 277 \cite{hillpeterson}. The claim developed here is
that force is naturally a covector, that impulse inherits that status, and that
Sutton's Eqs.~2--1 through 2--6 are the flat-space, fixed-direction,
constant-rate specialisation of a statement that does not need a metric.

\subsection{Work is a pairing, not a dot product}
Let $M$ be a smooth manifold, the configuration space of the mechanical system.
At each $x\in M$ the tangent space $T_xM$ holds velocities and the cotangent
space $T^*_xM$ holds linear functionals on them. The evaluation
$T^*_xM\times T_xM\to\R$, $(\alpha,v)\mapsto\alpha(v)$, is canonical: it is
defined by the manifold structure alone and requires \emph{no} metric.

A force is specified by the work it does along a path. For a curve
$C:[a,b]\to M$ with $t\mapsto x(t)$,
\[
 W=\int_C f=\int_a^b f\bigl(\dot x(t)\bigr)\,dt
  =\int_a^b f_k\bigl(x(t)\bigr)\frac{dx^k}{dt}\,dt ,
\]
which is the ordinary integral of the pullback $C^*f$ over $[a,b]$. The
integrand is a $1$-form $f=f_k\,dx^k\in\Omega^1(M)$, and the value is
independent of the parameterisation at fixed orientation. Frankel states the
conclusion directly: the integrand of a line integral is naturally a $1$-form,
not a vector, and the work integral ``does not require a metric''
\cite[Section O.f]{frankel}.

The familiar $W=\int\mathbf F\cdot d\mathbf r$ is therefore \emph{not} the
primitive statement. It presupposes a metric $g$ (to form the dot product) and
usually also Cartesian components (so that $g_{ij}=\delta_{ij}$ makes the
distinction invisible). Both are extra structure, and both are available in the
rocket problem --- which is precisely why the shortcut is safe there and unsafe
elsewhere.

\subsection{Why the index is down: the Legendre transform}
The covector status of force is not a convention; it follows from the
variational formulation. Take kinetic energy
$T=\tfrac12 m\,g_{ij}\dot x^i\dot x^j$ and a potential $V$ independent of
$\dot x$, with $L=T-V$.

\begin{proposition}\label{prop:momentum-covector}
The momentum conjugate to $x^k$ is
$p_k=\partial L/\partial\dot x^k=m\,g_{kj}\dot x^j$; that is,
$p=m\,v^\flat$, the metric dual of $m$ times the velocity.
\end{proposition}

\noindent The computation is immediate:
$\partial_{\dot x^k}\bigl(\tfrac12 m g_{ij}\dot x^i\dot x^j\bigr)
 =\tfrac12 m\bigl(g_{kj}\dot x^j+g_{ik}\dot x^i\bigr)=m g_{kj}\dot x^j$
by symmetry of $g$. Frankel performs exactly this calculation and concludes that
$p$ ``is $m$ times the covariant version of the velocity vector''
\cite[Section O.f]{frankel}.

The Legendre transform $\mathbb FL:TM\to T^*M$, $v\mapsto\partial L/\partial v$,
lands in the cotangent bundle by construction. The Euler--Lagrange equations,
augmented by a non-potential force,
\[
 \frac{d}{dt}\frac{\partial L}{\partial\dot x^k}
   -\frac{\partial L}{\partial x^k}=f_k ,
\]
place $f_k$ on the same footing as $\dot p_k$: both are components of elements
of $T^*_xM$. Momentum and force carry a lower index for the same structural
reason.

\subsection{Promotion and demotion: the musical isomorphisms}
A metric $g$ supplies --- and is exactly what supplies --- the identification
between the two spaces.

\begin{definition}
For a Riemannian metric $g$, define
$\flat:T_xM\to T^*_xM$ by $v^\flat=g(v,\cdot)$, so
$(v^\flat)_i=g_{ij}v^j$, and $\sharp:T^*_xM\to T_xM$ by
$(\alpha^\sharp)^i=g^{ij}\alpha_j$, where $g^{ij}g_{jk}=\delta^i_k$. The two are
mutually inverse.
\end{definition}

So one may always ``promote'' the force $1$-form to a vector field
$f^\sharp\in\mathfrak X(M)$, with $(f^\sharp)^i=g^{ij}f_j$, and Newton's second
law takes either the covector form $m(\nabla_{\dot x}\dot x)^\flat=f$ or the
vector form $m\nabla_{\dot x}\dot x=f^\sharp$. Nothing is lost --- but the
promotion is a choice that consumes the metric, and the two objects transform
differently under a change of coordinates.

In Cartesian coordinates on Euclidean $\R^3$, $g_{ij}=\delta_{ij}$ and the
components of $f$ and $f^\sharp$ are numerically equal. That coincidence is the
whole reason engineering mechanics can treat force as a vector without ever
being wrong. It stops being a coincidence in curvilinear coordinates. In
cylindrical coordinates $(r,\theta,z)$ on flat space,
\[
 g=\operatorname{diag}(1,\;r^2,\;1),\qquad
 f_\theta=r^2f^\theta,
\]
so the covariant, contravariant and \emph{physical} components
$f_{(\theta)}=f_\theta/r=r f^\theta$ are three different numbers. This is not
pedantry for turbomachinery: Hill and Peterson's angular-momentum balance,
$rF_\theta=\delta m\,\tfrac{d}{dt}(rc_\theta)$ \cite[Eq.~7.5]{hillpeterson},
carries exactly these $r$ weightings, and stating which component is meant is
what keeps a torque bookkeeping consistent.

\subsection{Impulse is the increment of a covector}
Write the impulse of $f$ over $[t_0,t_1]$ as
$J=\int_{t_0}^{t_1}f(t)\,dt$. There is an obstruction hiding in that
expression: $f(t)\in T^*_{x(t)}M$, and covectors at \emph{different} points lie
in different vector spaces, so the integrand cannot simply be summed. Making
$J$ well defined requires one of

\begin{itemize}\itemsep2pt
\item a flat space with a globally parallel (inertial Cartesian) frame, so all
      fibres are canonically identified; or
\item a connection, integrating the parallel-transported covector; or
\item contraction with a covariantly constant vector field first, reducing the
      integrand to a scalar.
\end{itemize}

\noindent Rocket performance uses the first, and the third is what the scalar
total impulse actually is.

\begin{proposition}\label{prop:impulse-momentum}
In Cartesian coordinates on Euclidean $\R^3$, $g$ is constant, hence
$\partial T/\partial x^k=0$ and the Euler--Lagrange equation reduces to
$\dot p_k=f_k$. Integrating,
$J=\int_{t_0}^{t_1}f\,dt=p(t_1)-p(t_0)$, an identity in $T^*\R^3$.
\end{proposition}

\noindent The hypothesis is doing real work, and it is worth seeing exactly
where the second term comes from, because it is easy to mistake it for a
product rule when it is really just Euler--Lagrange rearranged.

Since $V$ does not depend on $\dot x$, we have $p_k=\partial L/\partial\dot
x^k=\partial T/\partial\dot x^k$ and $\partial L/\partial x^k=\partial
T/\partial x^k-\partial V/\partial x^k$. Substituting both into the
Euler--Lagrange equation and absorbing the potential force into $f_k$,
\begin{equation}\label{eq:two-terms}
 \dot p_k=f_k+\frac{\partial T}{\partial x^k}.
\end{equation}
The extra term exists because $T=\tfrac12 m\,g_{ij}(x)\dot x^i\dot x^j$
depends on \emph{position}, through the metric. In Cartesian coordinates
$g_{ij}=\delta_{ij}$ is constant, $\partial T/\partial x^k=0$, and
\eqref{eq:two-terms} collapses to $\dot p_k=f_k$.

The product and chain rules appear one level down, in expanding the left side:
\[
 \dot p_k=\frac{d}{dt}\bigl(m\,g_{kj}(x)\dot x^j\bigr)
  =m\bigl(\partial_ig_{kj}\,\dot x^i\dot x^j+g_{kj}\ddot x^j\bigr),
\]
the first term by the product rule on $g_{kj}\dot x^j$ and then the chain rule
on $g_{kj}(x(t))$. With
$\partial T/\partial x^k=\tfrac12m\,\partial_kg_{ij}\dot x^i\dot x^j$,
equation~\eqref{eq:two-terms} becomes
\[
 m\,g_{kj}\ddot x^j
 +m\Bigl(\partial_ig_{kj}-\tfrac12\partial_kg_{ij}\Bigr)\dot x^i\dot x^j
 =f_k .
\]
Because $\dot x^i\dot x^j$ is symmetric in $i,j$, only the symmetric part of
$\partial_ig_{kj}$ survives, so $\partial_ig_{kj}\dot x^i\dot x^j
=\tfrac12(\partial_ig_{kj}+\partial_jg_{ki})\dot x^i\dot x^j$, and
\[
 m\,g_{kj}\ddot x^j+m\,\Gamma_{k,ij}\dot x^i\dot x^j=f_k,
 \quad
 \Gamma_{k,ij}=\tfrac12\bigl(\partial_ig_{kj}+\partial_jg_{ki}
   -\partial_kg_{ij}\bigr),
\]
the Christoffel symbol of the first kind. Raising the index recovers
$m(\ddot x^l+\Gamma^l_{ij}\dot x^i\dot x^j)=f^l$: the forced geodesic
equation, which is Newton's second law written covariantly.

\subsubsection*{The cylindrical case, explicitly}
Take $(r,\theta,z)$ with $g=\operatorname{diag}(1,r^2,1)$, so
$T=\tfrac12m(\dot r^2+r^2\dot\theta^2+\dot z^2)$. Then
\[
 p_r=m\dot r,\qquad p_\theta=mr^2\dot\theta,\qquad p_z=m\dot z,
\]
\[
 \frac{\partial T}{\partial r}=mr\dot\theta^2,\qquad
 \frac{\partial T}{\partial\theta}=0,\qquad
 \frac{\partial T}{\partial z}=0 .
\]
Equation~\eqref{eq:two-terms} in the radial slot gives
$m\ddot r=f_r+mr\dot\theta^2$, i.e.\ $f_r=m(\ddot r-r\dot\theta^2)$ --- the
centrifugal term is the $\partial T/\partial r$ contribution, not a fictitious
force added by hand. In the angular slot $\partial T/\partial\theta=0$, so
\[
 \dot p_\theta=f_\theta
 \;\Longleftrightarrow\;
 \frac{d}{dt}\bigl(mr^2\dot\theta\bigr)=f_\theta ,
\]
and expanding gives $m(r^2\ddot\theta+2r\dot r\dot\theta)=f_\theta$, whose
$2r\dot r\dot\theta$ is the Coriolis term.

Note what the lower index bought: $p_\theta=mr^2\dot\theta$ is the angular
momentum and $f_\theta$ is the torque, both with an extra factor of $r$
relative to the physical components $c_\theta=r\dot\theta$ and
$f_{(\theta)}=f_\theta/r$. Hill and Peterson's
$rF_\theta=\delta m\,\tfrac{d}{dt}(rc_\theta)$
\cite[Eq.~7.5]{hillpeterson} is exactly $\dot p_\theta=f_\theta$ written in
physical components, since $r\,c_\theta=r^2\dot\theta=p_\theta/m$ and their
$rF_\theta$ is our $f_\theta$. The covariant components are the ones an
engineer already calls angular momentum and torque; the index position is
telling you which quantity you are holding.

Because $\partial T/\partial x^k\neq0$ here, the naive componentwise
$\int f\,dt=\Delta p$ is false in these coordinates.

\subsection{From the covector to Sutton's total impulse}
Sutton defines total impulse by
$I_t=\int_0^t F\,dt$ \cite[Eq.~2--1]{sutton}, where $F$ is a scalar. The
geometric content is the third bullet above: if $\hat e$ is a fixed unit
direction (the thrust axis) and $F=f(\hat e)$, then
\[
 I_t=\int_0^{t}f(\hat e)\,dt=J(\hat e),
\]
one component of the impulse covector. Four assumptions separate this from the
general statement.

\begin{description}\itemsep2pt
\item[A1] Flat space, inertial Cartesian frame --- so $J$ exists
  (Prop.~\ref{prop:impulse-momentum}).
\item[A2] The thrust direction $\hat e$ is constant over $[0,t]$.
\item[A3] $F\ge0$, the thrust magnitude along that axis.
\item[A4] Start and stop transients are negligible, where Eq.~2--2 is used.
\end{description}

\begin{proposition}\label{prop:scalar-bound}
$\lVert J\rVert\le\int_0^t\lVert f\rVert\,dt$, with equality if and only if the
direction of $f$ is constant. Hence Sutton's scalar $I_t$ equals the magnitude
of the delivered vector impulse only under A2; for a gimbaled or attitude-slewing
burn it is a strict over-estimate.
\end{proposition}

\noindent That is the engineering content of insisting on the covector: a
quoted total impulse is a bound, not a delivered $\Delta v$ budget, unless the
thrust direction held still.

Sutton's Eq.~2--2, $I_t=Ft$, follows by additionally holding $F$ constant: with
$F$ constant the integral $\int_0^tF\,dt$ is just $Ft$, and A4 is what allows
the start and stop transients --- during which $F$ is emphatically not constant
--- to be left out of the integral rather than modelled. For specific impulse Sutton writes
\[
 I_s=\frac{\int_0^t F\,dt}{g_0\int_0^t\dot m\,dt}
 \qquad\text{\cite[Eq.~2--3]{sutton}},
\]
a ratio of two integrals. It is worth being explicit about what average that is.

\begin{proposition}\label{prop:weighted-mean}
With the instantaneous value $I_s(t)=F/(\dot m g_0)$,
\[
 \frac{\int F\,dt}{g_0\int\dot m\,dt}
 =\frac{\int I_s(t)\,\dot m\,dt}{\int \dot m\,dt},
\]
the \emph{mass-flow-weighted} mean of $I_s$, not its time average. The two
coincide only when $\dot m$ is constant.
\end{proposition}

\noindent Sutton's Eq.~2--5, $I_s=F/(\dot m g_0)$, is then the constant-$F$,
constant-$\dot m$ case, and Eq.~2--4, $I_s=I_t/(m_pg_0)$, follows from
$\int\dot m\,dt=m_p$.

Two further points of hygiene. First, $g_0$ is a \emph{defined} constant,
$9.80665\,\mathrm{m/s^2}$, not the local gravitational acceleration; Sutton says
so, and notes that in orbit $m_pg_0$ ``does not represent the weight''
\cite[Section 2.1]{sutton}. Consequently $I_s$ measured in seconds is not a
duration --- it is a velocity divided by a conventional constant, and the unit
is an artefact of that convention.

Second, and this is where the metric re-enters, Sutton's effective exhaust
velocity
\[
 c=I_sg_0=F/\dot m \qquad\text{\cite[Eq.~2--6]{sutton}}
\]
divides a force by a mass flow rate and calls the result a velocity. Force is a
covector and velocity is a vector, so the equality is really
$c\,\hat e=(f/\dot m)^\sharp$: the identification is the musical isomorphism,
performed silently. In Cartesian components on Euclidean space the two sides
have equal numbers, so nothing goes wrong. The statement to carry forward is
that $c$ is a metric-dependent object in a way that $I_t$, as a pairing, is not.

\subsection{Where the thrust covector comes from}
$f$ in the vehicle equation of motion is not primitive either: it is the
resultant of surface tractions --- pressure and viscous --- over the wetted
boundary. Traction is naturally a covector-valued $2$-form, the object Frankel
uses for stress and which pairs with a deformation $1$-form to give a scalar
$3$-form of work \cite[Section on Elasticity and Stresses]{frankel}. Sutton's
momentum-flux expression for thrust is a consequence of a control-volume
balance, not a definition of force, and the rocket's variable mass means the
vehicle-only $F=d(mv)/dt$ is \emph{not} a correct starting point. That
derivation belongs with Sutton Section 2.2 and is not attempted here.

\subsection{Notation: the symbol \texorpdfstring{$c$}{c}}
Reading Sutton and Hill and Peterson together creates a live collision.

\begin{description}\itemsep2pt
\item[Sutton] $c$ is the effective exhaust velocity \cite[Eq.~2--6]{sutton};
  $c^*$ is characteristic velocity. Sutton notes the Russian literature uses
  $c$ in place of $I_s$.
\item[Hill and Peterson] $c$ (and $\mathbf c$) is the \emph{absolute fluid
  velocity}, with components $c_r,c_\theta,c_z$ throughout the turbomachinery
  chapters; $c$ also denotes duct circumference, specific heat of a liquid and
  the speed of light. Speed of sound is $a$, and $c^*$ is again characteristic
  velocity \cite[Appendix X]{hillpeterson}.
\end{description}

\noindent So $c$ \emph{is} the propulsion convention for effective exhaust
velocity, and it is worth keeping rather than inventing a private symbol. The
collision is with Hill and Peterson's fluid velocity, not with the speed of
sound --- these notes already write $a$ for the latter, as in the nozzle topic.
The rule adopted here: $c$ denotes effective exhaust velocity in
rocket-performance contexts; subscripted $c_r,c_\theta,c_z$ denote absolute
fluid velocity components in turbomachinery contexts, following Hill and
Peterson; and where both appear in one derivation, the exhaust velocity is
written $v_e$ and said to equal $c$.

\subsection{Scope of what is claimed}
Propositions~\ref{prop:momentum-covector}--\ref{prop:weighted-mean} are
mathematical identities under their stated hypotheses, verifiable by direct
computation; they are not physical claims. A1--A4 are modelling assumptions
about a particular burn, and whether they hold is an engineering question about
the vehicle, not something the geometry decides. No numerical or symbolic study
accompanies this topic yet, so no executable evidence is claimed here; the
natural first check is Proposition~\ref{prop:scalar-bound} on a slewing burn,
comparing $\int\lVert f\rVert\,dt$ against $\lVert\int f\,dt\rVert$.

\subsection{An erratum, and why it is the relevant kind}
Sutton's worked Example~2--1 (printed p.~31 / PDF p.~55) prints
\[
 F=I_sg_0=23.3\times2354.4=54{,}936\ \mathrm{N},
\]
which was checked here against the source PDF's own text layer, not merely
against OCR. Two defects sit in that line. First, $I_sg_0=2354.4\
\mathrm{m/s}$ is the effective exhaust velocity $c$ --- the book computes it as
$c$ one sentence earlier --- so the equation sets a force equal to a velocity.
Eq.~2--6 rearranges to $F=\dot mc=\dot mI_sg_0$; the factor $\dot m$ is
missing. Second, the displayed multiplicand does not reproduce the stated
result: $23.3\times2354.4=54{,}857.5$, whereas the printed $54{,}936$ follows
from the unrounded $\dot m=70/3$, since $2354.4\times70/3=54{,}936$ exactly.

Every downstream quantity in the example is nevertheless correct, including
$I_t=164{,}808\ \mathrm{N\,s}$ --- independently confirmed by
$I_sm_pg_0=240\times70\times9.81$ --- and the final acceleration
$422.58\ \mathrm{m/s^2}=43.08\,g_0$. Only the symbolic statement is wrong, and
it survives into the ninth edition.

The erratum is worth recording because of \emph{what kind} of error it is: a
force and a velocity have been identified. That is exactly the confusion the
distinction between $f\in T^*M$ and $c\,\hat e=(f/\dot m)^\sharp\in TM$ is
there to prevent, and it is invisible to a reader who has already accepted that
force ``is a vector'' and that the metric may be applied silently. Carrying the
index position makes the line fail to typecheck on sight.
